% =============================================================================
% Refined \subsection{Qualitative Analysis of Divergent Cases}
% Uses real example content from assessment outputs as evidence.
% Paste this into your paper in place of the original subsection.
% =============================================================================

\subsection{Qualitative Analysis of Divergent Cases}

The qualitative cases illustrate why the relationship between solving and assessment is strong but incomplete. We consider the four outcome combinations (solve correct/incorrect $\times$ assess correct/incorrect), with particular attention to the two most informative types: \textbf{solve correct, assess incorrect} and \textbf{solve incorrect, assess correct}. We support each claim with concrete evidence from the model's outputs.

\paragraph{Solve correct, assess incorrect.}
In these cases the model produces the correct final answer when solving the problem itself but fails to identify the earliest erroneous step in the benchmark solution---so it detects that something is wrong but mislocalizes the error.

\textbf{GSM8K item gsm8k-192.} The problem asks: over 30 days, Christina had 12 good days initially; the first 24 days are given (8 good, 8 bad, 8 neutral); the next three days are good, neutral, good; how many good days were \emph{left} in the month? The benchmark solution concludes there are 0 good days left. The human annotator marks the \emph{first} error at step 2, where the solution states that ``the number of days left in the month is $30 - 24 = 6$ days'' and builds the rest of the reasoning on that. GPT-4, when assessing this solution, predicted the first error at step 5 instead. The model's rationale was: ``Step 5 incorrectly states that the next three days add 3 more good days to the count, but only two of the next three days are good (good, neutral, good), so it should add 2, not 3.'' Thus the model correctly identified a real counting error (adding 3 instead of 2 good days) but attributed it to step 5 rather than to the earlier step 2, where the structure of the reasoning first goes wrong. This shows that solving the problem correctly does not guarantee correct localization of the earliest error.

\textbf{MATH item math-509.} The problem is to compute $\tan 315^\circ$. The benchmark solution uses the reference angle $45^\circ$ and states that in the fourth quadrant ``the sine is positive and the cosine is negative,'' then gives $\tan 315^\circ = -1$. The human label marks the first error at step 2 (the quadrant sign claim: in the fourth quadrant, sine is negative and cosine is positive). GPT-4 predicted the first error at step 3 and gave the rationale: ``Step 3 contains a conceptual error: it incorrectly states that in the fourth quadrant, `the sine is positive and the cosine is negative.' In fact, in the fourth quadrant, sine is negative and cosine is positive.'' So the model identified the very same conceptual mistake as the annotator but pointed to step 3 instead of step 2. The assessment is wrong only in the \emph{index} of the earliest error, not in the type of error---illustrating that critique can be semantically on target while still failing the exact-match criterion used in the benchmark.

\paragraph{Solve incorrect, assess correct.}
Here the model fails to solve the original problem (its own solution has a wrong final answer) but correctly assesses the benchmark solution---either by correctly predicting ``no error'' (label $-1$) or by correctly identifying the earliest error step.

\textbf{GSM8K item gsm8k-238.} The problem asks: Charisma meditates 15 minutes twice per day and does 1 hour of yoga 5 days per week; how much time on meditation and yoga in 4 weeks? The benchmark solution computes $0.5$ hours/day meditation $\times$ 7 days $\times$ 4 weeks $= 14$ hours, plus 5 hours/week yoga $\times$ 4 weeks $= 20$ hours, for a total of 34 hours. All steps are correct (human label $-1$). GPT-4's \emph{own} solution to this problem was incorrect (solve incorrect), but when given the benchmark solution to assess, the model predicted label $-1$ and stated: ``All steps are correct. Each calculation accurately reflects the problem conditions, and the reasoning is valid throughout.'' So the model failed to generate a correct solution itself yet correctly judged this solution as error-free, indicating that assessment can succeed using capabilities that are not sufficient for correct generative solving.

\textbf{MATH item math-572.} The problem gives $a+b=7$ and $a^3+b^3=42$ and asks for $1/a + 1/b$ as a common fraction. The benchmark solution uses the sum-of-cubes identity, derives $ab = 43/3$, then computes $1/a + 1/b = (a+b)/(ab) = 7/(43/3) = 21/43$. The human label is $-1$ (all steps correct). GPT-4 failed to solve this problem in the solving task but, when assessing the given solution, predicted $-1$ with the rationale: ``All steps are mathematically correct: the sum of cubes identity is applied properly, the expressions are manipulated correctly, and the final answer is computed accurately using the given values.'' Again, the model's generative solution was wrong while its step-level assessment was correct, supporting the view that critique can draw on skills that partially overlap with, but are not identical to, those required for solving.

\paragraph{Summary.}
The four cases, with the exact problem content, step indices, and model rationales above, show that (i) solving correct does not imply assessment correct (gsm8k-192, math-509: wrong earliest-error localization despite correct final answer), and (ii) solving incorrect does not imply assessment incorrect (gsm8k-238, math-572: correct ``no error'' assessment despite wrong own solution). Together they support the claim that transfer from problem-solving to step-level assessment is real but incomplete, and that assessment relies in part on capabilities that are not fully redundant with generative solving.
